% Part: first-order-logic
% Chapter: introduction
% Section: semantic-notions

\documentclass[../../../include/open-logic-section]{subfiles}

\begin{document}

\olfileid{fol}{int}{sem}

\olsection{चिन्हार्थविषयक संकल्पना}

$\Sat{M}{!A}[s]$ लागू आहे का याचा विचार करताना, सोयीसाठी आपण $s$ कडून
सर्व !!{variable}s ना मूल्ये नेमून घेतो, पण~$!A$ मधील !!{variable}s ना
त्याने नेमून दिलेली मूल्येच फक्त वापरली जातात, असे आपण वर म्हटले.
खरे तर, $!A$ मधील \emph{मुक्त} चरांची मूल्येच महत्त्वाची असतात. आपण
काटेकोर असल्यामुळे अर्थातच हे तथ्य सिद्ध करणार आहोत. !!{sentence}s मध्ये
मुक्त चरे नसल्यामुळे, !!a{structure} मध्ये त्यांची पूर्ति होते की नाही
यासाठी $s$ चा काहीही फरक पडत नाही. म्हणून, $!A$ हे !!a{sentence} असेल,
तेव्हा ``प्रत्येक~$s$ साठी $\Sat{M}{!A}[s]$'' असा $\Sat{M}{!A}$ चा अर्थ
परिभाषित करता येतो; आणि तो योगायोगाने किमान एका~$s$ साठी
$\Sat{M}{!A}[s]$ असते तेव्हा आणि केवळ तेव्हाच सत्य असतो. !!{formula}s
साठी पूर्तिची कार्यक्षम व्याख्या मिळवण्यासाठी आपल्याला !!{variable} ची
मूल्यांकने आणावी लागतात, पण !!{sentence}s साठी पूर्ति ही !!{variable}
च्या मूल्यांकनांपासून स्वतंत्र असते.

``$\Sat{M}{!A}$'' ची व्याख्या मिळाल्यानंतर प्रथम-क्रम तर्कशास्त्रातील
!!{sentence}s च्या संदर्भात ``प्रसंग'' आणि ``यात सत्य'' यांचा अर्थ काय
हे आपल्याला माहीत होते. मग !!{sentence}s साठीच्या $\Sat{M}{!A}$ च्या
व्याख्येच्या आधारे वैधता, तार्किक निष्पन्नता आणि पूर्ततायोग्यता या मूलभूत
चिन्हार्थविषयक संकल्पना परिभाषित करता येतात. प्रत्येक !!{structure}
एखाद्या वाक्याची पूर्ति करत असेल, तर ते वाक्य वैध, $\Entails !A$, असते.
~$\Gamma$ मधील सर्व !!{sentence}s ची पूर्ति करणारे प्रत्येक
!!{structure} हे~$!A$ चीही पूर्ति करत असेल, तर !!{sentence}s च्या संचातून
ते तार्किकरीत्या निष्पन्न होते, $\Gamma \Entails !A$. आणि एखादे
!!{structure} एखाद्या !!{sentence}s च्या संचातील सर्व !!{sentence}s ची एकाच
वेळी पूर्ति करत असेल, तर तो संच पूर्ततायोग्य असतो.

!!{formula}s विगमनाधारित रीतीने परिभाषित केली आहेत आणि पूर्ति हीही
!!{formula}s च्या संरचनेवरील विगमनाने परिभाषित केली आहे; म्हणून आपल्या
चिन्हार्थमीमांसेचे गुणधर्म सिद्ध करण्यासाठी आणि परिभाषित
चिन्हार्थविषयक संकल्पनांचा परस्परसंबंध लावण्यासाठी विगमन वापरता येते.
यांपैकी काही गुणधर्म आपण एकत्र करून सिद्ध करू; काही अंशी प्रत्येक
गुणधर्म स्वतंत्रपणे रोचक असल्यामुळे, पण मुख्यतः चिन्हार्थमीमांसा आणि
!!{derivation} पद्धती यांच्यातील संबंधाचा अभ्यास पुढे करताना त्यांपैकी
अनेक उपयोगी पडणार असल्यामुळे. हे करण्यासाठी आपण अजून उल्लेख न केलेल्या
काही विन्यासविषयक संकल्पना आणि क्रियाही (काटेकोरपणे, म्हणजे विगमनाने)
परिभाषित कराव्या लागतील.

\end{document}
